% Source: Solutions/Yash-Indrajit-Himanshu-PS6.pdf, pp. 2-3 (typed, joint submission; Q4-6 credited to Indrajeet Kumawat).
\submittedby{Indrajeet Kumawat}

\begin{proof}
A surjective continuous map $f\colon X\to Y$ is a quotient map if for any subset
$U\subseteq Y$, $U$ is open in $Y$ if and only if $f^{-1}(U)$ is open in $X$ (equivalently with
closed sets). We analyze the two cases separately:

\textbf{Case 1 ($f$ is an open map):} Let $U\subseteq Y$ be any subset. If $U$ is open in $Y$,
then $f^{-1}(U)$ is open in $X$ by continuity. Conversely, suppose $f^{-1}(U)$ is open in $X$.
Since $f$ is an open map, the image $f(f^{-1}(U))$ is open in $Y$. Because $f$ is surjective,
$f(f^{-1}(U))=U$, so $U$ is open in $Y$.

\textbf{Case 2 ($f$ is a closed map):} Let $C\subseteq Y$ be any subset. If $C$ is closed in
$Y$, then $f^{-1}(C)$ is closed in $X$ by continuity. Conversely, suppose $f^{-1}(C)$ is closed
in $X$. Since $f$ is a closed map, $f(f^{-1}(C))$ is closed in $Y$. Since $f$ is surjective,
$f(f^{-1}(C))=C$, so $C$ is closed in $Y$.

\textbf{Conclusion:} In both cases, the quotient map condition holds. Therefore, $f$ is a
quotient map.
\end{proof}
